% !PRINTED 577
\weberpage{577}
Man kann die Form dieser Entwickelung aus den Produktformeln [\S\ 24, (11)]:
\[
\lambda = q^{-\frac{1}{2}} \Pi\left(\frac{1+q^{2\nu-1}}{1+q^{2\nu}}\right)^4 + 16\,q^{\frac{1}{2}}\ \Pi\left(\frac{1+q^{2\nu}}{1+q^{2\nu-1}}\right)^4
\]
ableiten, und findet:
\weq{3}{\lambda = q^{-\frac{1}{2}}(1 + \lambda_1 q + \lambda_2 q^2 + \lambda_3 q^3 + \cdots),}
worin die Konstanten $\lambda_1, \lambda_2, \lambda_3, \ldots$ \gesperrt{ganze rationale Zahlen} sind.

Setzen wir ferner
\weq{4}{x = \sqrt{\varkappa}\ \operatorname{sn}\left(\frac{w}{\sqrt{\varkappa}},\varkappa\right),}
so bleibt (4) nach \S\ 45, (9) ungeändert, wenn $\varkappa$ mit $1:\varkappa$ vertauscht wird. Es genügt $x$ der Differentialgleichung
\weq{5}{\left(\frac{dx}{dw}\right)^2 = 1 - \tfrac{1}{4}\lambda x^2 + x^4.}
Entwickelt man also $x$ nach dem \gesperrt{Taylorschen} Lehrsatz in eine Reihe nach steigenden Potenzen von $w$:
\weq{6}{x = w + X_1 w^3 + X_2 w^5 \ldots,}
so sind die Koeffizienten $X_1, X_2, \ldots$ \gesperrt{ganze rationale Funktionen von $\lambda$ mit rationalen Zahlenkoeffizienten}, wie sich aus (5) nach der Methode der unbestimmten Koeffizienten ergibt.

Benutzt man die Bezeichnung von \S\ 42, so ist
\weq{7}{x = \sqrt{\varkappa}\ \operatorname{sn} v = \frac{\vartheta_{11}(u)}{\vartheta_{01}(u)},}
und nach \S\ 57, IV. mit einer etwas modifizierten Bedeutung der $A, B, C, D$:
\weq{8}{\begin{aligned}
\sqrt{\varkappa}\ \operatorname{sn} mv &= \frac{xA(x)}{D(x)}, \\
\operatorname{cn} mv &= \sqrt{1-\frac{x^2}{\varkappa}}\ \frac{B(x)}{D(x)}, \\
\operatorname{dn} mv &= \sqrt{1-\varkappa x^2}\ \frac{C(x)}{D(x)}.
\end{aligned}}
Hier sind $A(x), B(x), C(x), D(x)$ ganze rationale Funktionen von $x$ vom Grade $\frac{1}{2}(m^2-1)$, und es ist aus den Rekursionsformeln \S\ 57, (13), (14) zu ersehen, daß die Funktionen $A, D$ und das Produkt $BC$ ganze rationale Funktionen von $\lambda$ sind, deren
% [source sic: the printed degree (m^2-1)/2 conflicts with the m^2-1 roots and later degree formulae]
%CONT>>
